\documentclass[12pt]{article}
\usepackage[utf8]{inputenc}
\usepackage{amsmath}
\usepackage{amssymb}
\usepackage{hyperref}
\usepackage{setspace}
\onehalfspacing

% A numbered claim and the section that argues it (check_claims.py).
\newcommand{\claim}[2]{\item #2 (Section~\ref{#1}).}

\title{Unity requires physically enforced agreement}
\author{Sebastian Bobadilla-Suarez\\
Independent Researcher\\
\texttt{sebastian.bobadilla.s@gmail.com}\\
ORCID: \href{https://orcid.org/0000-0002-5951-0772}{0000-0002-5951-0772}}
\date{}

\begin{document}
\maketitle

\begin{abstract}
Human conscious experience is unified: the brain's many local descriptions of a
shared situation agree. We argue that unity requires this agreement to be
physically enforced, produced and maintained by a coupling through which the
graded physical state of each region continuously pulls the regions it overlaps
toward agreement. Clocked digital hardware computes agreement instead. Its
noise margins make every content readout locally constant in micro state, and
we prove that a readout with a margin, under continuous micro dynamics, depends
on no region's graded state; content read from bits is therefore never
physically enforced, whatever algorithm the bits run. For a coupled sheet of
oscillators we show that the width of the coupling kernel sets how fast
agreement decays with distance, and that only a kernel tail slower than
$r^{-4}$ makes it independent of distance. The premise yields interventional
tests on cortex and on artificial systems, and we state the outcomes that would
refute it on each.
\end{abstract}

\section{Introduction}
\label{sec:intro}

\paragraph{Premise P (physical unity).} A system's contents are unified only if
the content of each region depends continuously on the graded physical state of
the regions it overlaps, through a coupling that contracts their disagreement.

P is a necessary condition for unity. It rejects functionalism about unity: a
perfect simulation of a coupled field reproduces the field's content-level
dynamics and fails P. The paper makes three claims:
\begin{enumerate}
\claim{sec:enforcement}{A content readout with a noise margin, under continuous micro dynamics, depends on the graded state of no region, so content read from bits is not physically enforced whatever algorithm the bits run}
\claim{sec:reach}{On a noisy coupled sheet, the width of the coupling kernel sets how fast agreement decays with distance, and only a kernel tail slower than $r^{-4}$ makes agreement independent of distance}
\claim{sec:tests}{P predicts that a sub-threshold perturbation of one cortical region shifts the content decoded in an overlapping region toward agreement, and that no perturbation below a digital system's noise margin shifts any of its contents}
\end{enumerate}

\section{Synchrony is not unity}
\label{sec:synchrony}

Equal phases coexist with incompatible local descriptions, and compatible local
descriptions determine a unique global state. Unity is therefore a property of
what overlapping regions describe, and synchrony is at most a route to it.

\section{Unity as agreement over a declared cover}
\label{sec:cover}

Agreement is measured between regions that decode a shared variable. The cover
of regions is fixed by decodability before agreement is measured, so that the
measurement cannot be emptied by the choice of cover.

\section{Physical enforcement versus routed agreement}
\label{sec:enforcement}

A readout has a noise margin on an operating set when every micro state in that
set has a neighbourhood on which the readout is constant. Content depends
gradedly on a region when no neighbourhood of that region's micro state leaves
the next content unchanged. If the micro dynamics is continuous and carries a
state into a margin's operating set, no region's micro state enters the next
content gradedly (\texttt{not\_physicallyEnforced\_of\_margin}). A thresholded
continuous quantity has a margin off its threshold, and a finite word of such
readouts has one too, so the conclusion holds for any content read from bits
(\texttt{not\_physicallyEnforced\_of\_bits}).

\section{How far a physical coupling enforces agreement}
\label{sec:reach}

At identical natural frequencies a noisy coupled sheet is a gradient system. In
the harmonic approximation the variance of the phase difference between two
sites is the noise strength times their effective resistance, with the coupling
constants as conductances. Under short-range coupling that resistance grows as
the logarithm of distance, with a coefficient inversely proportional to the
kernel's second moment; a kernel tail slower than $r^{-4}$ bounds it.

\section{The unity window}
\label{sec:window}

Agreement cannot arrive before the causal cone of the coupling allows, and
disagreement decays at the relaxation rate of the coupling. Unity at a given
time scale requires the interval between the two to be nonempty.

\section{Current artificial systems}
\label{sec:ai}

A transformer running on clocked digital hardware reads every content from bits
held within their noise margins, so the result of
Section~\ref{sec:enforcement} applies to it. Hardware whose contents are graded
physical quantities coupled without barriers, such as analog, in-memory or
neuromorphic substrates, can satisfy P.

\section{Tests}
\label{sec:tests}

In cortex, weak stimulation of one region below the behavioural threshold
should shift the content decoded in an overlapping region, gradedly and toward
agreement; an unchanged decoded content in the overlapping region counts
against P. In an artificial system, activation patching below and above the
scale at which the logical trajectory changes separates perturbations inside
the margin from those that cross it.

\section{Limitations}
\label{sec:limitations}

P is a premise and is not derived from the dynamics; the formal results show
what follows from it and not that it holds. It is necessary for unity and not
sufficient for consciousness: coupled pendulums satisfy it. The margin theorem
concerns idealised digital readouts, and graded leakage in real hardware that
carries content-relevant agreement would narrow its reach. The kernel result
rests on the harmonic approximation at identical natural frequencies and is
established numerically on finite sheets, not proved. Human data can falsify P
but cannot by themselves confirm it against functionalism, because graded
coupling and information flow co-vary in every manipulation available so far.

\end{document}
